% =============================================================================
% PREFACIO
% =============================================================================

\chapter*{Prefacio}
\addcontentsline{toc}{chapter}{Prefacio}

% --- Cita de apertura ---
\begin{flushright}
\itshape
``Bad programmers worry about the code. Good programmers worry about data structures and their relationships.''\\[4pt]
\upshape --- Linus Torvalds
\end{flushright}
\bigskip

% --- Texto introductorio ---
Este compendio de notas surge del estudio sistemático de las estructuras de datos
fundamentales del curso de \textit{Data Structures} (CISE-Ph.D.), con el objetivo de
documentar no solo los algoritmos y las interfaces de programación, sino la lógica
estructural que subyace a cada decisión de diseño. Para un investigador con formación
en física y métodos computacionales, las estructuras de datos no son colecciones
arbitrarias de convenciones: son soluciones a problemas de organización de información
bajo restricciones de tiempo y espacio, análogas en espíritu a las estructuras variacionales
o los esquemas numéricos que minimizan error bajo restricciones de recursos.

El hilo conductor de este texto es la pregunta: \textit{¿cuándo elegir qué estructura, y
a qué costo?} Cada capítulo responde esa pregunta para una familia de estructuras,
articulando la complejidad asintótica, los casos de uso y los compromisos de diseño de
forma precisa.

% --- Subsección: Audiencia objetivo ---
\subsection*{El investigador al que está dirigido este compendio}

Este documento está diseñado para estudiantes que se introducen en CISE y científicos
computacionales con madurez matemática propia de la física teórica y experimental. Se
asume familiaridad con análisis asintótico básico y programación orientada a objetos en
Java, pero no experiencia previa con las estructuras aquí cubiertas. El nivel de formalidad
es deliberadamente alto: se incluyen definiciones recursivas, notación de conjuntos y
argumentos de complejidad, no como adorno, sino porque son la única forma precisa de
caracterizar el comportamiento de estas estructuras.

% --- Subsección: Alcance y estructura ---
\subsection*{Alcance del compendio y estructura de las notas}

Este compendio documenta una fracción del temario total del curso, que abarca desde las
estructuras lineales elementales hasta las estructuras jerárquicas y los mecanismos de
búsqueda eficiente. El temario completo incluye, entre otras estructuras:

\begin{itemize}
    \item \textbf{Estructuras lineales:} arrays, matrices, listas enlazadas
          (\textit{linked lists}), pilas (\textit{stacks}), colas (\textit{queues}),
          \textit{deques} y conjuntos (\textit{sets}).
    \item \textbf{Estructuras asociativas:} tablas hash (\textit{hash tables}),
          mapas (\textit{hash maps}) y su implementación en Java.
    \item \textbf{Estructuras jerárquicas:} árboles binarios (\textit{binary trees})
          y árboles binarios de búsqueda (\textit{binary search trees}, BST).
\end{itemize}

Los capítulos aquí incluidos cubren los bloques asociativos y jerárquicos. A medida que
el curso avanza, el compendio se irá extendiendo para documentar el resto del temario
de forma coherente y acumulativa.

% --- Subsección: Anglicismos técnicos ---
\subsection*{Sobre los anglicismos técnicos}

En el ámbito de las estructuras de datos y los algoritmos, términos como \textit{hashing},
\textit{heap}, \textit{stack}, \textit{BST}, \textit{pre-order} o \textit{in-order} poseen
definiciones estandarizadas en la literatura internacional que no tienen equivalentes en
español suficientemente precisos o de uso extendido. A lo largo de este texto se conservan
en inglés o se usan como préstamos integrados, señalados en cursiva en su primera aparición.

% --- Subsección: Convenciones notacionales ---
\subsection*{Convenciones notacionales}

\begin{itemize}
    \item \textbf{Complejidad asintótica:} se emplea notación Big-O ($O(\cdot)$),
          Big-Theta ($\Theta(\cdot)$) y Big-Omega ($\Omega(\cdot)$) para caracterizar
          complejidad temporal y espacial en el peor caso, caso promedio y mejor caso,
          respectivamente.
    \item \textbf{Árboles:} $h(v)$ denota la altura del nodo $v$; $d(v)$ su profundidad.
          $T_v$ denota el subárbol con raíz en $v$.
    \item \textbf{Conjuntos:} $k(v)$ denota la clave almacenada en el nodo $v$.
          Los subárboles izquierdo y derecho se denotan $v.\mathit{left}$ y
          $v.\mathit{right}$ para mantener coherencia con la notación de implementación.
    \item \textbf{Pseudocódigo:} los algoritmos se presentan en pseudocódigo estructurado
          independiente del lenguaje, con implementaciones de referencia en Java donde
          resulta ilustrativo.
\end{itemize}

\vspace{1.2cm}
\begin{flushright}
\textit{Jesús Antonio Chala Casanova}\\[2pt]
CISE, \the\year
\end{flushright}
